\chapter{Vulvodynia}

\section*{Definition}
Vulvodynia is chronic vulvar pain lasting at least 3 months without an identifiable cause (infectious, dermatologic, neoplastic, or neurologic). It may be generalized (entire vulva) or localized (typically the vestibule, termed provoked vestibulodynia), and is diagnosed when other causes are excluded.

\section*{What Happens in Your Body}
Vulvodynia involves neuropathic pain mechanisms, with increased density of nociceptive nerve fibers (C-fibers) in the vestibular mucosa, peripheral and central sensitization, and elevated inflammatory immune signaling proteins (TNF-alpha, IL-6). Pelvic floor muscle hypertonicity and psychosocial factors (anxiety, catastrophizing) amplify and maintain the pain cycle.

\section*{Causes}
\begin{itemize}
  \item \textbf{Primary:} unknown cause; no identifiable organic cause
  \item \textbf{Neurologic:} Pudendal neuralgia, nerve entrapment, complex regional pain syndrome, peripheral neuropathy
  \item \textbf{Inflammatory:} Elevated local mast cells and proinflammatory immune signaling proteins (IL-6, TNF-alpha)
  \item \textbf{Musculoskeletal:} Pelvic floor hypertonicity, myofascial trigger points, pubic symphysis dysfunction
  \item \textbf{Hormonal:} Oral contraceptive use (altered hormone sensitivity of the vestibule), menopausal estrogen loss
  \item \textbf{Psychological risk factors:} Anxiety, depression, pain catastrophizing, history of abuse, sexual dysfunction
  \item \textbf{Genetic:} IL-1 receptor antagonist gene polymorphisms may increase susceptibility
\end{itemize}

\section*{When to See a Doctor}
See your doctor if:
\begin{itemize}
  \item Vulvar pain has persisted for 3 months or longer without clear cause
  \item Pain interferes with sexual activity, tampon use, sitting, or exercise
  \item There is associated burning, stinging, or raw sensation that is provoked or spontaneous
\end{itemize}

\section*{Related Symptoms}
\begin{itemize}
  \item pain during intercourse, vaginismus, pelvic floor hypertonicity, chronic pelvic pain, interstitial cystitis, fibromyalgia, irritable bowel syndrome, temporomandibular joint disorder
\end{itemize}
\textbf{Important:} Vulvodynia affects up to 10--16\% of women at some point in their lives, yet most do not seek medical care.

\section*{Self-Care / Home Management}
\begin{itemize}
  \item Avoid irritants: soaps, douches, scented products, tight synthetic underwear
  \item Apply cool gel packs to the vulva after provocation
  \item Use a water-based lubricant without glycerin or parabens
  \item Wear loose, cotton underwear and avoid prolonged sitting
\end{itemize}

\section*{How Your Doctor Will Evaluate You}
\begin{itemize}
  \item History: pain characteristics, provoking factors, associated conditions, past treatments
  \item Q-tip test: gentle palpation of the vestibule in clock-face distribution to map allodynia
  \item Rule out other causes: wet mount and culture (Candida, bacterial vaginosis, trichomonas), biopsy if lichen sclerosus or neoplasia suspected
  \item Assessment of pelvic floor muscle tone and trigger points
\end{itemize}

\section*{Treatment Options}
\begin{itemize}
  \item First-line: Multimodal approach including pelvic floor physical therapy and topical lidocaine
  \item Topical medications: compounded estradiol/testosterone (vestibulodynia), amitriptyline gel, gabapentin cream
  \item Oral medications: tricyclic antidepressants (amitriptyline, nortriptyline), gabapentin, pregabalin
  \item Vestibulectomy for refractory localized provoked vestibulodynia (surgical excision of the hymen and posterior vestibule)
  \item Psychological: CBT for pain management, biofeedback, mindfulness-based pain reduction
  \item Trigger point injections with lidocaine and corticosteroid
\end{itemize}

\section*{References}
\begin{thebibliography}{99}
\bibitem{issvd-vulvodynia} Bornstein J, et al. "ISSVD terminology and classification of vulvodynia." \emph{J Low Genit Tract Dis}. 2024;28(1):1-10.
\bibitem{uptodate-vulvodynia} Haefner HK, et al. Vulvodynia: clinical features and diagnosis. In: UpToDate, Post TW (Ed), UpToDate, Waltham, MA; 2025.
\bibitem{nejm-vulvodynia} Pukall CF, et al. "Vulvodynia." \emph{N Engl J Med}. 2023;389(23):2180-2190.
\bibitem{bmj-vulvodynia} Andrews JC, et al. "Vulvodynia." \emph{BMJ}. 2024;386:e081765.
\bibitem{cochrane-vulvodynia} Andrews JC, et al. "Interventions for vulvodynia." \emph{Cochrane Database Syst Rev}. 2023;(10):CD015238.
\bibitem{obgyn-vulvodynia} Sadownik LA, et al. "Provoked vestibulodynia: evidence-based treatment." \emph{Obstet Gynecol}. 2024;143(2):231-242.
\end{thebibliography}
\clearpage
